% SCOPE CONTRACT: PERSON A (Empirical Case)
\section{Empirical Case: Onocoy}
\label{sec:onocoy_case}

\subsection{System Overview}
Onocoy operates a decentralized GNSS correction network based on Real-Time Kinematic (RTK) positioning. RTK combines reference-station correction streams with rover-side measurements to reduce positioning error from meter-scale to centimeter-scale ranges, which is relevant for surveying, precision agriculture, robotics, and industrial automation use cases.
% TODO-CITE: [Teunissen and Montenbruck 2017] RTK precision and GNSS correction fundamentals

The network architecture depends on geographically distributed reference stations rather than a single centralized infrastructure operator. This topology is designed to improve coverage and redundancy by sourcing correction data from many independent stations connected through a common protocol layer.
% TODO-CITE: [Onocoy docs / whitepaper] distributed station architecture and service model

\subsection{Network Participants}
Onocoy includes two primary participant groups. First, reference-station operators (often labeled miners in DePIN contexts) deploy and maintain GNSS station hardware, maintain connectivity, and provide reliable uptime/quality signals. Station operation is not purely passive: installation quality, calibration, and maintenance discipline affect delivered correction quality.
% TODO-CITE: [Onocoy docs; RTCM standards references] operator requirements and quality constraints

Second, data users (rover-side consumers) purchase correction access for high-accuracy positioning workflows. Typical demand comes from professional/industrial use cases where service continuity and measurement quality are operational requirements rather than optional features.
% TODO-CITE: [Onocoy docs; EUSPA GNSS market report] user categories and demand profile

\subsection{Tokenomic Mechanics (Descriptive)}
Onocoy uses a dual-layer design with ONO as the transferable utility/governance token and Data Credits as non-transferable utility credits for service consumption. Data Credits are purchased with fiat-denominated value, then consumed and burned during service usage.
% TODO-CITE: [Onocoy whitepaper Rev 3.0.1, token architecture sections] dual-token and Data Credit mechanics

Supply-side incentive issuance follows a decay schedule for ONO. The whitepaper describes a capped ONO supply with a deflationary release profile (16\% annual deflation factor, equivalent to a four-year halving profile for newly created units), while buyback-and-burn behavior is tied to protocol-level revenue routing choices.
% TODO-CITE: [Onocoy whitepaper Rev 3.0.1] cap, release schedule, and buyback/burn routing

On public documentation reviewed so far, staking/collateral rules for miner participation are not documented with the same precision as emission and data-credit mechanics. This remains an open descriptive item to confirm from current protocol docs before final freeze.
% TODO-CITE: [Onocoy docs staking/miner participation requirements] confirm current staking rule set

\subsection{Person A Onocoy Draft Import (Provisional, Source-Preserving)}
\label{subsec:personA_onocoy_import}
% Provisional source-preserving import from Person A draft snapshot.
% Source: tmp/docs/personA_submission_extracted.txt
% Expanded integration pass to maximize contributor visibility before final synthesis.

The following imported block preserves expanded descriptive material from Person A's Onocoy case draft.

Onocoy is described as a decentralized Real-Time Kinematic (RTK) GNSS correction network. RTK correction is used to improve positioning precision from meter-level to centimeter-level ranges, with relevance for precision agriculture, robotics and automation, surveying, and industrial machinery control.
% TODO-CITE: [Person A sources: Teunissen and Montenbruck 2017; EUSPA 2023]

The draft emphasizes that Onocoy coordinates distributed reference stations operated by independent participants. These stations provide correction data that is aggregated and distributed through a protocol-mediated coordination layer.

The demand side is characterized as primarily professional and enterprise-oriented. In this view, demand for high-accuracy correction services is linked to operational workflows and continuity requirements rather than to purely speculative token narratives.
% TODO-CITE: [Person A sources: EUSPA 2023; Onocoy docs]

Provider requirements are also preserved explicitly. Reference-station operators are described as responsible for:
\begin{itemize}
    \item deploying compatible GNSS hardware,
    \item maintaining continuous uptime,
    \item maintaining calibration quality,
    \item providing stable connectivity.
\end{itemize}

Failure to maintain uptime or quality can reduce delivered service reliability and affect reward eligibility.
% TODO-CITE: [Person A sources: RTCM 2023; Onocoy docs]

Person A's token architecture framing is preserved as a dual-layer model:
\begin{itemize}
    \item ONO token for utility/governance and incentive-layer functions.
    \item Data Credits as non-transferable usage credits consumed during service access.
\end{itemize}

Users purchase Data Credits for service usage, and credits are consumed on use. This architecture is presented as an attempt to separate user-side payment flows from direct exposure to secondary-market token volatility.
% TODO-CITE: [Person A sources: Onocoy docs; Helium docs comparison]

The draft also notes that Onocoy remains exposed to broader crypto-market cycles and macro conditions despite utility-demand anchoring. Token volatility can still influence provider-side incentives and retention behavior.
% TODO-CITE: [Person A sources: Cong et al. 2021]

A further preserved analytical point is that Onocoy's central stress-relevant interaction is between relatively stable industrial demand and potentially unstable token-market dynamics. This interaction is treated as a key bridge to later stress-model parameterization and scenario testing.

\paragraph{Traceability note.} Full verbatim Person A submission snapshots are preserved in Appendix~\ref{subsec:contributor_draft_archive} for audit and attribution during drafting.

